\section{讨论}\label{sec:discussion}

\subsection{三轴不是一条权衡曲线}

题目把问题写成“可解释性 vs. 性能”，隐含了一个单维权衡：解释性换来性能，或者反过来。
本文的实验不支持这个图景。在最短路这个任务上，白盒在三个轴上\emph{同时}占优：

\begin{table}[H]
\centering
\caption{三轴上的对照（各行的图集不同，末列注明 $n$；L1 为 128 轮，L2 为 8 层）}
\label{tab:three}
\small
\setlength{\tabcolsep}{4pt}
\begin{tabular}{llllc}
\toprule
轴 & L0 / L1（白盒） & L2（黑盒） & 口径 & $n$ \\
\midrule
准确率 & $\sim10^{-6}$ / $0.024$ & 11.77（中位 14.92，最好 2.19） & 同分布 ER & 40 \\
鲁棒性 & $\sim10^{-6}$ / $0.035$ & 22.93 $\to$ 22.83 & 权重噪声 20\% & 30 \\
可解释性 & 精度 1.00，成路率 100\% & 精度 0.50，成路率 11\% & 同预算 $k$ & 28 \\
\bottomrule
\end{tabular}
\end{table}

这张表能说与不能说要分清。\textbf{能说}：在“允许结构臂迭代到收敛（128 轮）”这个设定下，白盒在三轴上占优。
\textbf{不能说}：白盒在任何设定下都占优——表~\ref{tab:same} 已经给出反例：同深度（8 层）时，结构臂在 SBM 上反而更差（35.38 对 28.76，$p<0.001$），在网格与规模外推上不显著（$p=0.095$）。
换句话说，白盒的胜势有一半来自“它本来就不需要层数这个约束”，而不是“结构比黑盒更会学”。

关于计算开销，本文没有做公平的对比，因此不放进表里：白盒真正的成本是跑一次 Dijkstra（$n=800$ 时 4.1\,ms，见第~\ref{sec:white}~节），而 L0 之所以要跑 128 轮，是因为我们\emph{把它写成了模型}；真实场景里没有人会那样用白盒。

所以正确的问题不是“要解释还是要性能”，而是：\textbf{你知不知道算法}。知道，就把它写进结构里，三轴一起赢；不知道，才需要在解释与性能之间做取舍。

\subsection{可验证性缺口}

本文建议用一个更锋利的量来替代“可解释性”这个模糊的词：\textbf{误差有没有上界}。

白盒的误差是零，而且扰动之后的新答案落在可计算的区间里（命题~\ref{prop:mono}--\ref{prop:bound}）。
黑盒的误差没有上界：在本文的设定里，它可以输出负距离——把真值约 57 的一对点算成 $-1.9$——而没有任何机制阻止它这么做。
模型越“黑”，我们能对它的输出说的话就越少——这不是精度问题，是\emph{可验证性}问题。

这也解释了为什么 L1 那个 0.15\% 的参数偏差值得单独讨论：在一个迭代算法里，“几乎正确”不是一个稳定的状态，误差会被迭代放大。
白盒的正确性是\emph{结构性}的（恒等映射写死在结构里），黑盒的正确性是\emph{近似性}的，两者在单步上可能看不出差别，但在 128 步之后就分开了。

\subsection{结论的边界}\label{sec:limits}

这一节的每一条都是对本文结论的削弱，但必须写。

\begin{enumerate}[leftmargin=1.8em]
\item \textbf{基线没有调到最好。}L2 的 6/8 个种子塌缩成常数预测，说明这个 mean-pool 基线训练得并不好。它最好的一次是 1.90，仍比 L1 差两个数量级，但“GNN 学不了最短路”这个更强的结论，本文的数据\emph{不支持}。能支持的只是：在这个未调好的基线上，结构先验带来 2--3 个数量级的差距。
\item \textbf{参数量不对称。}L1 有 2 个参数，L2 有 $104{,}321$ 个。这是刻意的：结构对齐的意义就是“用先验换样本”，硬把 L1 扩到同等参数量等于人为削弱它。但读者应当把参数量摆在心里看表。
\item \textbf{迭代预算不对称。}L1 在推理时跑 128 轮（迭代到收敛），L2 只有训练时的 8 层。这不是黑盒的固有属性——如果把 L2 的每层参数共享，它也可以跑 128 轮——而是我们的设计选择。因此 L1 与 L2 的比较是联合比较，不是受控消融。
\item \textbf{任务是被刻意选的。}最短路有精确解、有已知算法、真值可精确计算。这恰好是白盒的主场。换成“给定部分观测预测用户会走哪条路”这类任务，白盒的先验就不再完整，权衡曲线会重新出现。
\item \textbf{可解释性指标本身有陷阱。}本文用了一个反例说明 Fidelity$^+$ 会被常数预测器骗过。这说明“可解释性”没有公认的单一度量，任何单一数字都可能被刷。我们用三个与任务绑定的指标来交叉验证，但这仍不是终局。
\item \textbf{只跑了一组超参。}L1 的训练只用了 40 轮、学习率 $10^{-2}$、固定初始化，没有做超参扫描，也没有画 $(a,b)$ 的损失面。因此“训练停在 0.9988 而不是 1”这个观察，不能排除“换一组设置就能收敛”。我们保留这个观察，但不把它上升为“学习必然失败”的结论。
\item \textbf{有一篇直接的反例必须提。}Nerem 等人~\cite{nerem2025ood} 证明：带合适正则的 GNN 可以\emph{精确}实现 Bellman-Ford，并在规模上外推。这与本文 L0 的结论一致（结构对了就精确），但同时说明本文 L2 的失败更像“基线没调好”，而不是“GNN 学不会最短路”。本文的强结论只到“未调好的 mean-pool 基线做不到”为止。
\end{enumerate}

\subsection{什么时候黑盒才值得}

按本文的逻辑，黑盒的价值应当出现在“算法未知或不可写”的地方：
目标函数没有多项式算法、输入带有感知噪声（图像、文本、传感器）、或者人自己也说不清规则、只能给例子。
在这些场合，白盒根本没有可写的东西，讨论“误差有没有界”也就无从谈起——此时黑盒是唯一选择，可解释性则退化为一种\emph{审计}工具：不是去证明它对，而是去发现它在什么情况下会错。

反过来，如果一个任务已经有精确算法，而模型的作用只是加速或近似它，那么最划算的做法不是“用 GNN 学这个算法”，而是“把这个算法写进结构里，让学习只去补它缺的那部分”。
本文的 L0 就是前者的极端情形：当结构已经完整时，学习不但没有增益，还会破坏它。
